\section{Threat-Campaign Scenarios}\label{sec:threat-campaigns}

Threat campaigns combine delivery, user interaction, carrier structure,
execution, and payload behavior. A file-defense component observes only a
bounded segment of that chain. The scenarios below are therefore analytical
compositions for corpus construction and reason-family reporting. They are not
claims that \AFD{} attributes an attachment to a named actor or infers a malware
family. ATT\&CK techniques provide a vocabulary for initial access,
masquerading, and user execution~\cite{mitreT1566001,mitreT1036002,
mitreT1036007,mitreT1036008,mitreT1204002}. The enforceable predicates come
from the local attachment.

\subsection{Masqueraded active attachment}

In the first scenario, an operator sends an executable, script, shortcut, or
active document under a media-themed display name. The name may use a double
extension or a right-to-left override so that the rendered suffix appears
benign. The supplied media type may claim an image while byte evidence identifies
an executable or unsupported grammar. The campaign may seek user execution of a
Trojan, downloader, dropper, stager, ransomware program, or spyware implant.
Those downstream labels and roles do not change the local decision.

The observable defense is deterministic name validation and concrete-profile
convergence. Path separators, control and bidirectional override characters,
dangerous nonfinal suffixes, and forbidden final suffixes violate the evaluated
name policy. A conflict among normalized name, supplied media type, and bytes
is assigned a blocking decision before publication. If evaluation records that
decision with no published bytes, the permitted conclusion is that the
masqueraded or unsupported carrier did not become a deliverable artifact at
this boundary. If the same payload arrives through a channel outside \AFD, or
if the endpoint is already compromised, this conclusion does not apply.

\subsection{Staged delivery roles}

Staged campaigns require three separate operational labels. A downloader
obtains follow-on material from another source after execution. A dropper
carries follow-on material inside itself and places that embedded payload after
execution. A stager prepares an execution context or transfers control to a
follow-on stage. A stager may download or drop its next stage, but downloader,
dropper, and stager are not interchangeable labels.

Archives, scripts, shortcut files, installers, and active office documents can
carry any of these roles because they can encode commands or package multiple
objects. The media-only release profile does not attempt to sanitize these
classes. It rejects them as unsupported. This is a deliberate scope choice
because preserving their intended semantics would preserve the active behavior
that makes them useful to the campaign.

The permitted claim is containment at one attachment boundary, not observation
of a staged role. \AFD{} does not observe network retrieval, extraction of an
embedded payload, command execution, or transfer to a later stage. A benign
unsupported archive is blocked by the same rule. The evaluation must therefore
measure policy enforcement and false blocking relative to the frozen media-only
policy rather than infer downloader, dropper, or stager intent from rejection.

\subsection{Exploit-bearing media}

An exploit-bearing media scenario uses a nominally supported image, animation,
audio stream, or video container to reach a vulnerable parser or decoder.
Mutated dimensions, lengths, offsets, checksums, sample tables, lacing, frame
counts, and nested elements can test structural paths. Validly encoded codec
content can also exercise a decoder vulnerability without violating the outer
grammar. This distinction determines the claim.

Bounded input parsing, checked arithmetic, reconstruction, independent output
parsing, process timeouts, output-size limits, and clean-only publication reduce
specific opportunities for malformed-container delivery and unbounded work.
A semantic image or GIF path can emit new encoded content from bounded decoded
state. Structural audio and video paths retain codec payloads and support only
container and metadata claims unless policy requires semantic reconstruction.
If the required semantic path is unavailable or fails, semantic-or-block policy
rejects the artifact. A sample-specific success or rejection does not prove
that arbitrary accepted bytes are harmless, and a decoder invocation remains
inside the trusted computing and fault-containment assumptions.

\subsection{True-polyglot and boundary-smuggling delivery}

A true-polyglot campaign exploits a whole octet string that is completely
accepted by at least two incompatible parsers. This differs from boundary
smuggling. A valid media object with an executable, archive, or inert foreign
object after its terminal boundary is an appended or trailing-object carrier,
not a true dual-format polyglot, unless both parsers accept the whole string
completely. An executable prefix before a later media signature is prefix
camouflage under the same rule. Prior work shows both detector-side parsing
differences and semantic gaps among implementations~\cite{jana2012,koch2025,
you2025}.

\AFD{} requires parsing from offset zero and complete consumption through the
exact terminal boundary for prefix and suffix constructions. It treats at
least two incompatible complete parses of the whole string as true-polyglot
ambiguity. Location-aware allowlists govern nested structures. A structural
handler cannot transform a retained ambiguous region into a neutralized claim.
A semantic path may support the term ``neutralized'' only when no independently
valid active grammar is authorized, no source container or compressed payload
bytes are copied, and the new output has one canonical interpretation under the
registered validators. Otherwise the correct action and report are
``blocked.''

\subsection{Payload-objective strata}

Campaign corpora often organize samples by the final payload, including
ransomware or spyware. This organization remains useful for coverage of threat
provenance, but it must be crossed with carrier class. A ransomware executable
masquerading as a PNG exercises active-class and classifier-disagreement rules.
An image with location metadata exercises privacy-leakage metadata exclusion
only if that field is part of the frozen profile. A spyware label attached to
the same sample remains provenance context and does not convert passive
metadata into spyware. A ransomware-labelled exploit inside a valid codec
stream tests the exact semantic or structural path used by that format. The
payload label does not alter what evidence the pipeline collects.

\begingroup
\footnotesize
\setlength{\tabcolsep}{4pt}
\begin{longtable}{L{0.15\textwidth}L{0.20\textwidth}L{0.31\textwidth}L{0.23\textwidth}}
\caption{Threat-campaign compositions and attachment-level claim boundaries.}
\label{tab:campaign-scenarios}\\
\toprule
Scenario & Local observables & Enforced response under the evaluated profile & Residual or excluded claim \\
\midrule
\endfirsthead
\toprule
Scenario & Local observables & Enforced response under the evaluated profile & Residual or excluded claim \\
\midrule
\endhead
\midrule
\multicolumn{4}{r}{Continued on next page}\\
\endfoot
\bottomrule
\endlastfoot
Masqueraded active file & Bidi control, dangerous suffix chain, false media type, executable or unsupported bytes & Reject name, profile disagreement, or unsupported grammar and expose no deliverable output & No attribution and no malware-label or operational-role classification \\
Downloader carrier & Script, shortcut, archive, package, installer, active document, or executable & Reject unsupported active class & No observation that the component fetched follow-on material \\
Dropper carrier & Script, shortcut, archive, package, installer, active document, or executable & Reject unsupported active class & No observation that the component contained and placed an embedded payload \\
Stager carrier & Script, shortcut, archive, package, installer, active document, or executable & Reject unsupported active class & No observation that the component prepared or transferred control to a later stage \\
Malformed media exploit & Invalid bounds, offsets, checksums, order, terminator, or over-budget structure & Block parse or reconstruction failure, or publish only independently validated reconstruction & No universal codec-vulnerability coverage \\
Valid codec exploit & Supported outer grammar with decoder-triggering content & Require qualified semantic reconstruction for a bitstream-disarm claim, otherwise block & Structural copying does not neutralize retained codec bytes \\
True polyglot & Whole octet string completely accepted by at least two incompatible parsers & Reject incompatible complete parses, or require qualified semantic reconstruction with one accepted output grammar & Retained pixels and samples may preserve information \\
Boundary-smuggling carrier & Prefix camouflage, appended or trailing object, or nested active element & Enforce zero-offset parsing, exact consumption, and location-aware allowlists & The construction is not thereby a true dual-format polyglot \\
Privacy-bearing media & Excluded EXIF, XMP, comments, application identifiers, or other named metadata & Rebuild without fields forbidden by the concrete output profile & Passive leakage is not spyware and visible or steganographic sensitive content may remain \\
\end{longtable}
\endgroup

\subsection{Evaluation use and nonclaims}

Each scenario becomes a set of manifest records with independent fields for
delivery context, behavioral ground truth, carrier class, concrete format,
mutation recipe, expected action, and expected reason family. Crossed cases are
necessary. The same carrier mutation should appear with benign and malicious
provenance where licensing and safety controls permit, and the same payload
stratum should appear across multiple carrier forms. This design tests whether
the verdict follows observable policy predicates rather than a convenient
directory label.

The campaign analysis yields no empirical effectiveness number in the drafting
environment. Tests, coverage-guided fuzzing, fault injection, corpus execution,
fidelity analysis, differential validation, and performance measurements are
reserved for the pinned server run. Until that run is complete, the tables in
this section specify expected mechanisms and claim boundaries only. After the
run, reporting must preserve separate counts for unsupported-class blocking,
classifier disagreement, malformed structure, resource rejection, polyglot
ambiguity, prefix camouflage, trailing-object rejection, nested-object policy,
qualified semantic neutralization, privacy-metadata exclusion, and benign false
blocking. Combining these into one malware-blocking percentage would erase the
causal distinction that the taxonomy is designed to preserve.
